% =====================================================================
%  Quantum part - Exercise 1 : selecting a heralding scheme
% =====================================================================

\subsection{Exercise 1 -- Selecting a heralding scheme}
\label{sec:q1}

\paragraph{Notation and common quantities.}
The group velocity in the fiber is $v = c/n = \SI{2e8}{\metre\per\second}$, i.e.\ \SI{5}{\micro\second\per\kilo\metre}.
In MIM each photon travels $L/2$ to the midpoint BSA. The efficiency of one arm is the product of the
local efficiency (extraction/coupling and detection) and of the fiber transmission:
\begin{equation}
  \eta = \eta_\mathrm{loc}\,T(L/2),
  \qquad \eta_\mathrm{loc} = \eta_\mathrm{ext}\eta_\mathrm{cpl}\eta_\mathrm{det} = 0.5\times0.8 = 0.4,
  \qquad T(\ell) = 10^{-\alpha \ell/10}.
\end{equation}
\begin{center}
\begin{tabular}{lccc}
  \toprule
  $L$ & $\alpha L/2$ & $T(L/2)$ & $\eta = \eta_\mathrm{loc}\,T(L/2)$ \\
  \midrule
  \SI{40}{\kilo\metre}  & \SI{4}{dB}  & 0.398 & 0.159 \\
  \SI{100}{\kilo\metre} & \SI{10}{dB} & 0.100 & 0.040 \\
  \bottomrule
\end{tabular}
\end{center}

% ---------------------------------------------------------------------
\paragraph{1. Effective MIM attempt frequency.}
With one memory per endpoint, a new attempt cannot start before the BSA outcome is received.
The photon travels $L/2$ to the midpoint and the classical result travels $L/2$ back, so the
MIM feedback delay is
\begin{equation}
  t_\mathrm{link} = \frac{L/2 + L/2}{v} = \frac{nL}{c}
  = \begin{cases}
      \dfrac{1.5\times\SI{40e3}{\metre}}{\SI{3e8}{\metre\per\second}} = \SI{200}{\micro\second} & (L=\SI{40}{\kilo\metre})\\[2ex]
      \SI{500}{\micro\second} & (L=\SI{100}{\kilo\metre})
    \end{cases}
\end{equation}
The entire cost of an attempt is the local trial-clock period plus the feedback delay, so
$f_\mathrm{eff,MIM} = 1/(t_\mathrm{link} + t_\mathrm{clock})$:
\begin{equation}
  f_\mathrm{eff}(\SI{40}{\kilo\metre}) = \frac{1}{\SI{210}{\micro\second}} = \SI{4.76}{\kilo\hertz},
  \qquad
  f_\mathrm{eff}(\SI{100}{\kilo\metre}) = \frac{1}{\SI{510}{\micro\second}} = \SI{1.96}{\kilo\hertz}.
\end{equation}

% ---------------------------------------------------------------------
\paragraph{2. Dual-click MIM.}
Both memories must prepare a memory--photon pair (probability $q$ each), both photons must be
detected (probability $\eta$ each), and the BSA accepts the outcome with probability $p_\mathrm{BSA}$:
\begin{equation}
  p_\mathrm{dual} = q^2\, p_\mathrm{BSA}\, \eta^2 = 0.32\,\eta^2,
  \qquad R_\mathrm{dual} = p_\mathrm{dual}\, f_\mathrm{eff}.
\end{equation}
\begin{align}
  p_\mathrm{dual}(\SI{40}{\kilo\metre}) &= 0.8^2\times0.5\times0.159^2 = \num{8.11e-3},
  & R_\mathrm{dual}(\SI{40}{\kilo\metre}) &= \SI{38.6}{\per\second},\\
  p_\mathrm{dual}(\SI{100}{\kilo\metre}) &= 0.8^2\times0.5\times0.04^2 = \num{5.12e-4},
  & R_\mathrm{dual}(\SI{100}{\kilo\metre}) &= \SI{1.0}{\per\second},
\end{align}
obtained as $\num{8.11e-3}\times\SI{4.76}{\kilo\hertz}$ and $\num{5.12e-4}\times\SI{1.96}{\kilo\hertz}$ respectively.

% ---------------------------------------------------------------------
\paragraph{3. Single-click MIM: largest excitation probability $p_1$.}
Each memory emits a photon with probability $p_1$. A single click at the midpoint is accepted.
\begin{itemize}
  \item \emph{True herald}: exactly one memory emits and its photon is detected,
        $p_\mathrm{true} = 2\eta\,p_1(1-p_1)$.
  \item \emph{False herald}: both memories emit, but only one of the two photons is detected
        (a single click is still observed), $p_\mathrm{false} = 2\eta(1-\eta)\,p_1^2$.
\end{itemize}
The accepted-click probability is $p_\mathrm{click} = p_\mathrm{true} + p_\mathrm{false}$ and the
conditional Bell-state fidelity is
\begin{equation}
  F_\mathrm{Bell} = \frac{p_\mathrm{true}}{p_\mathrm{true}+p_\mathrm{false}}
  = \frac{2\eta p_1(1-p_1)}{2\eta p_1(1-p_1) + 2\eta(1-\eta)p_1^2}
  = \frac{1-p_1}{1-p_1\eta}.
\end{equation}
Imposing $F_\mathrm{Bell} \ge F_\mathrm{min}$ and solving for $p_1$:
\begin{equation}
  p_1 \le \frac{1-F_\mathrm{min}}{1-F_\mathrm{min}\,\eta} = \frac{0.02}{1-0.98\,\eta}
  \;\Rightarrow\;
  p_1(\SI{40}{\kilo\metre}) = 0.0237, \qquad p_1(\SI{100}{\kilo\metre}) = 0.0208 .
\end{equation}

% ---------------------------------------------------------------------
\paragraph{4. Single-click accepted heralded-state rate.}
\begin{equation}
  p_\mathrm{click} = 2\eta(1-p_1)p_1 + 2\eta(1-\eta)p_1^2,
  \qquad R_\mathrm{single} = p_\mathrm{click}\, f_\mathrm{eff}.
\end{equation}
\begin{center}
\begin{tabular}{lccccc}
  \toprule
  $L$ & $p_\mathrm{true}$ & $p_\mathrm{false}$ & $p_\mathrm{click}$ & $F_\mathrm{Bell}$ & $R_\mathrm{single}$ \\
  \midrule
  \SI{40}{\kilo\metre}  & \num{7.37e-3} & \num{1.50e-4} & \num{7.52e-3} & 0.980 & \SI{35.8}{\per\second} \\
  \SI{100}{\kilo\metre} & \num{1.63e-3} & \num{3.33e-5} & \num{1.66e-3} & 0.980 & \SI{3.26}{\per\second} \\
  \bottomrule
\end{tabular}
\end{center}

% ---------------------------------------------------------------------
\paragraph{5. Preferred implementation.}
\begin{center}
\begin{tabular}{lccc}
  \toprule
  $L$ & Dual-click & Single-click & Ratio single/dual \\
  \midrule
  \SI{40}{\kilo\metre}  & \textbf{\SI{38.6}{\per\second}} & \SI{35.8}{\per\second} & 0.93 \\
  \SI{100}{\kilo\metre} & \SI{1.0}{\per\second} & \textbf{\SI{3.26}{\per\second}} & 3.3 \\
  \bottomrule
\end{tabular}
\end{center}
At \SI{40}{\kilo\metre} dual-click is preferred: the attempt frequency is the same for both schemes, and
dual-click heralds states with ideal fidelity, whereas single-click must keep $p_1$ small and still
accepts a $2\%$ fraction of false heralds.
At \SI{100}{\kilo\metre} single-click is preferred: $p_\mathrm{dual}\propto\eta^2$ while
$p_\mathrm{click}\propto\eta$, so when the channel becomes very lossy the dual-click probability
collapses much faster. The single-click rate becomes about 3 times larger, despite the fidelity
constraint on $p_1$.

% ---------------------------------------------------------------------
\paragraph{6. Crossover distance.}
Both schemes are MIM, so they share the same $f_\mathrm{eff}$ and $R_\mathrm{single} = R_\mathrm{dual}$
reduces to $p_\mathrm{click} = p_\mathrm{dual}$. Substituting $p_1 = 0.02/(1-0.98\eta)$:
\begin{equation}
  p_\mathrm{click} = 2\eta p_1(1-p_1\eta) = \frac{0.04\,\eta(1-\eta)}{(1-0.98\eta)^2}
  \quad\text{and}\quad
  \frac{0.04\,\eta(1-\eta)}{(1-0.98\eta)^2} = 0.32\,\eta^2
  \;\Rightarrow\;
  8\eta(1-0.98\eta)^2 = 1-\eta .
\end{equation}
Solving numerically gives $\eta = 0.145$. Then
\begin{equation}
  T(L/2) = \frac{0.145}{\eta_\mathrm{loc}} = \frac{0.145}{0.4} = 0.363
  \;\Rightarrow\; 10^{-0.02\,L/2} = 0.363
  \;\Rightarrow\; L \approx \SI{44}{\kilo\metre}.
\end{equation}
Below about \SI{44}{\kilo\metre} dual-click gives the higher rate; above it, single-click does.

% ---------------------------------------------------------------------
\subsubsection*{Single-click MM extension ($L = \SI{40}{\kilo\metre}$)}
The interference station is moved to endpoint B. The photon from A traverses the whole link, while
the photon from B follows only the local optical path:
\begin{equation}
  \eta_A = \eta_\mathrm{loc}\,T(L) = 0.4\times10^{-0.2\times40/10} = 0.0634,
  \qquad
  \eta_B = \eta_\mathrm{loc} = 0.4,
  \qquad
  \frac{\eta_A}{\eta_B} = T(L) = 0.158 .
\end{equation}

\paragraph{7. Balanced excitation probabilities.}
With $p_{0,i} = 1-p_{1,i}$, the two detected true-herald contributions are
$p_{1,A}\,p_{0,B}\,\eta_A$ (only A emits and its photon is detected) and $p_{0,A}\,p_{1,B}\,\eta_B$.
Balancing them:
\begin{equation}
  p_{1,A}(1-p_{1,B})\,\eta_A = p_{1,B}(1-p_{1,A})\,\eta_B
  \quad\Leftrightarrow\quad
  \frac{p_{1,B}}{1-p_{1,B}} = 0.158\,\frac{p_{1,A}}{1-p_{1,A}} .
\end{equation}
Under this condition
\begin{equation}
  p_\mathrm{true} = 2\,p_{1,A}\,\eta_A(1-p_{1,B}),
  \qquad
  p_\mathrm{false} = p_{1,A}\,p_{1,B}\,(\eta_A+\eta_B-2\eta_A\eta_B),
\end{equation}
where the false herald requires both emissions and exactly one of the two photons detected.
The fidelity constraint $F_\mathrm{Bell} = p_\mathrm{true}/(p_\mathrm{true}+p_\mathrm{false}) \ge F_\mathrm{min}$
is equivalent to
\begin{equation}
  \frac{p_\mathrm{false}}{p_\mathrm{true}}
  = \frac{p_{1,B}(\eta_A+\eta_B-2\eta_A\eta_B)}{2\eta_A(1-p_{1,B})}
  \le \gamma = \frac{1-F_\mathrm{min}}{F_\mathrm{min}} = 0.0204 .
\end{equation}
Solving for $p_{1,B}$, and then using the balance condition for $p_{1,A}$:
\begin{align}
  p_{1,B} &= \frac{2\gamma\eta_A}{\eta_A+\eta_B-2\eta_A\eta_B+2\gamma\eta_A} = 0.00623,\\
  p_{1,A} &= \frac{p_{1,B}\,\eta_B}{\eta_A + p_{1,B}(\eta_B-\eta_A)} = 0.0381 .
\end{align}
As expected $p_{1,A} > p_{1,B}$: the photon from A suffers the full-link loss $T(L)$, so A
must be excited about six times more often than B for the two contributions to reach the detector with
equal weight.

\paragraph{8. Herald probabilities and fidelity.}
\begin{align}
  p_\mathrm{true} &= p_{1,A}p_{0,B}\eta_A + p_{1,B}p_{0,A}\eta_B
     = \num{2.398e-3} + \num{2.398e-3} = \num{4.795e-3} \quad(\text{balanced}),\\
  p_\mathrm{false} &= p_{1,A}p_{1,B}(\eta_A+\eta_B-2\eta_A\eta_B) = \num{9.79e-5},\\
  p_\mathrm{click} &= p_\mathrm{true}+p_\mathrm{false} = \num{4.893e-3},\\
  F_\mathrm{Bell} &= \frac{\num{4.795e-3}}{\num{4.893e-3}} = 0.980 \;\ge F_\mathrm{min}\ \checkmark
\end{align}

\paragraph{9. MM attempt frequency and rate.}
In MM the photon from A travels the full length $L$ to the BSA at B, and the result must travel $L$
back to A before A can retry:
\begin{equation}
  2t_\mathrm{link} = \frac{n\,2L}{c} = \frac{1.5\times\SI{80e3}{\metre}}{\SI{3e8}{\metre\per\second}} = \SI{400}{\micro\second},
  \qquad
  f_\mathrm{eff,MM} = \frac{1}{2t_\mathrm{link}+t_\mathrm{clock}} = \frac{1}{\SI{410}{\micro\second}} = \SI{2.44}{\kilo\hertz},
\end{equation}
\begin{equation}
  R_\mathrm{MM} = p_\mathrm{click}\,f_\mathrm{eff} = \num{4.893e-3}\times\SI{2.44}{\kilo\hertz} = \SI{11.9}{\per\second}.
\end{equation}

\paragraph{10. Comparison at \SI{40}{\kilo\metre}.}
\begin{center}
\begin{tabular}{lccc}
  \toprule
  Scheme & Heralding delay & $f_\mathrm{eff}$ & Accepted rate \\
  \midrule
  MIM dual-click        & \SI{200}{\micro\second} & \SI{4.76}{\kilo\hertz} & \SI{38.6}{\per\second} \\
  MIM single-click      & \SI{200}{\micro\second} & \SI{4.76}{\kilo\hertz} & \SI{35.8}{\per\second} \\
  MM single-click (balanced) & \SI{400}{\micro\second} & \SI{2.44}{\kilo\hertz} & \SI{11.9}{\per\second} \\
  \bottomrule
\end{tabular}
\end{center}
\emph{Feedback time.} In MIM the station is at the midpoint, so photon and result travel
$L/2 + L/2$ (\SI{200}{\micro\second}), while in MM they travel $L + L$ (\SI{400}{\micro\second}).
With one memory per node no new attempt can start before the outcome is known, so the attempt
frequency drops from \SI{4.76}{\kilo\hertz} to \SI{2.44}{\kilo\hertz}.

\emph{Preparation balancing.} To equalise the true-herald weights, node A (whose photons suffer the
loss of the whole link) must emit much more often than node B ($p_{1,A}=0.0381$, $p_{1,B}=0.0062$).
This keeps the false heralds under control, but the probability per attempt is limited by the long,
lossy arm: $p_\mathrm{click,MM} = \num{4.89e-3}$, lower than both $p_\mathrm{click}=\num{7.52e-3}$ of
single-click MIM and $p_\mathrm{dual} = \num{8.11e-3}$ of dual-click MIM.

Losing both in probability per attempt and in attempt frequency, balanced single-click MM reaches only
about $1/3$ of the MIM rates. Dual-click MIM remains the best choice at \SI{40}{\kilo\metre}, thanks to
the midpoint station, which halves the feedback time, and to its heralded states being free of false heralds.
